\section{The anatomy layer: exact unwinding}\label{sec:anatomy}

It remains to reinstate the cofactor-anatomy indicator $\sigma$ in
$V_p(\lambda)$. By the conditional architecture, $\sigma$ here is the
\emph{anatomy-only} part of the cofactor weight: within-side digit
conditions belong to Lemma~C's uniformity clause (this interface is part
of the statements of \S\ref{sec:architecture}). The governing structural
fact makes the friable-sums literature unnecessary for the dominant
configuration.

\begin{proposition}[at most one large prime]\label{prop:one-large}
\statusPM\quad
Let $a\le A=x^{1-u}$ be a cofactor in the dominant configuration and
$y=(2x)^{1/3}$ the band floor. Since $u\ge\tfrac13$ implies $y^2\ge A$,
every such $a$ has at most one prime factor exceeding $y$, counted with
multiplicity. Consequently the friability indicator unwinds
\emph{exactly}:
\[
\mathbf 1\big[P^{+}(a)\le y\big]
\;=\;1-\sum_{\substack{\ell>y\ \mathrm{prime}\\ \ell\mid a}}1
\qquad(a\le A).
\]
\end{proposition}

\begin{proof}
If $\ell_1\ell_2\mid a$ with $\ell_1,\ell_2>y$ (allowing
$\ell_1=\ell_2$), then $a\ge y^2\ge A$, a contradiction (the boundary
case forces $a=y^2$, excluded as $y^2>A$ for $u>1/3$ and handled
trivially at $u=1/3$). So the sum on the right is $0$ or $1$, and equals
$1$ exactly when $P^{+}(a)>y$.
\end{proof}

\begin{remark}
In Dickman terms the cofactor's friability parameter is
$u_f=\log A/\log y=3(1-u)\in(1.5,2)$: strictly below $2$, which is why
the inclusion--exclusion truncates at the first term. The identity was
verified with zero violations on the actual progressions
(Appendix~\ref{app:numerics}).
\end{remark}

\begin{proposition}[decomposition into classical species]\label{prop:species}
\statusPM\quad
For any phase $\theta$, with $t_\ell$ the least index $t$ with
$\ell\mid a_0+qt$ (a single class mod $\ell$, since $(q,\ell)=1$),
\[
W(\theta):=\sum_{t\le R}\mathbf 1\big[P^{+}(a_0+qt)\le y\big]\,\e{t\theta}
\;=\;G(\theta)\;-\;\sum_{y<\ell\le A}\e{t_\ell\theta}
\sum_{s\le(R-t_\ell)/\ell}\e{s\,\ell\theta},
\]
where $G(\theta)=\sum_{t\le R}\e{t\theta}$. The identity is exact
(machine-checked to $10^{-14}$). Consequently D$^\dagger$-anatomy splits
by the dyadic size $\Lambda$ of $\ell$ into:
\begin{itemize}
\item\textbf{small/medium $\Lambda$} (many terms per $\ell$): geometric
sums at the \emph{dilated} frequencies $\ell\theta_\mu$ --- the same
Vinogradov-type counting as \S\ref{sec:minor}; the resonant $\ell$
(those with $\nm{\ell\theta_\mu}$ small) form a third, \emph{linear},
major-arc family, to be counted with Families A--B;
\item\textbf{large $\Lambda$} (one term per $\ell$): writing
$a=\ell b$, Type-II bilinear sums over $(\ell,b)$ with
$\ell b\equiv a_0\ (\mathrm{mod}\ q)$ --- standard bilinear
cancellation; at the extreme $\Lambda\sim A$ ($b$ bounded): primes in
arithmetic progressions mod $q$ against linear phases
$\e{b\theta_\mu s}$ --- Vinogradov/Balog technology, with $q=x^{u'}$
comfortably inside its range. Transition zones via the usual Vaughan
splitting.
\end{itemize}
\end{proposition}

\begin{proof}
By Proposition~\ref{prop:one-large},
$\mathbf 1[P^+(a)\le y]=1-\sum_{\ell>y,\,\ell\mid a}1$ for every $a$ in
range; insert into $W(\theta)$ and exchange sums; for fixed $\ell$ the
solutions of $\ell\mid a_0+qt$, $t\le R$, are $t=t_\ell+s\ell$. The
species classification is bookkeeping on the resulting double sum.
\end{proof}

\begin{remark}[status and fallback]
\statusN\quad
What remains for D$^\dagger$-anatomy is classical-toolbox write-up:
(i) the Type-I/II ranges with uniformity in $(q,\mu,p)$ and the Vaughan
gluing; (ii) the $\ell$-resonance count (third major-arc family, linear,
hence easier than Family A); (iii) the one genuine check: Lemma~B's
finer anatomy classes (sizes and counts of primes in sub-bands)
condition beyond bare friability --- each class adds
inclusion--exclusion layers of the \emph{same species}, but the
per-class verification must be done. Fallback if (iii) misbehaves: the
friable-exponential-sum literature, already pinned --- Harper
\cite{Harper16} (minor/major dichotomy for smooth Weyl sums), de la
Bret\`eche--Tenenbaum \cite{dlBT07} (friable sums at rational
arguments), Fouvry--Tenenbaum \cite{FT91} (friable integers in
progressions), Drappeau \cite{Drap15}, Drappeau--Shao \cite{DS16}, de la
Bret\`eche--Granville \cite{dlBG22}. Note our cofactor condition (``no
prime in the band above $p$'' plus size constraints) is smooth-\emph{type},
not literally $y$-friable: an adaptation layer, not a citation, even in
the fallback.
\end{remark}
